\documentclass[10pt,a4paper]{amsart}
\usepackage[margin=19mm]{geometry}
\usepackage{amsmath,amssymb,mathtools}
\usepackage[scr=rsfs,cal=euler]{mathalfa}
\usepackage{xfrac}
\usepackage[unicode,hidelinks,bookmarks=false]{hyperref}
\newcommand{\ie}{i.e.\ }
\newcommand{\cf}{cf.\ }
\newcommand{\resp}{respectively\ }
\newcommand{\Subsection}{\subsection}
\providecommand{\nobreakdash}{\nobreak\hbox{-}}
\setlength{\emergencystretch}{2em}
\multlinegap0pt
\usepackage{amssymb}
\usepackage{bm}
\usepackage{enumerate}
\usepackage{xcolor}
\usepackage[normalem]{ulem}
\newtheorem{thm}{Theorem}
\title{Limits of Stochastic Semigroups and Block-Triangular Majorisation}
\author{Fabio Deelan Cunden, Jakub Czartowski, Giovanni Gramegna, Marilena Ligab\`o}
\date{Source: arXiv:2609.04057v1 (CC BY 4.0)}
\begin{document}
\maketitle

\noindent\emph{Source and licence.} Limits of Stochastic Semigroups and Block-Triangular Majorisation. Version arXiv:2609.04057v1 (CC BY 4.0), distributed by arXiv under the Creative Commons Attribution 4.0 International licence, http://creativecommons.org/licenses/by/4.0/, as recorded in the arXiv licence metadata for this exact version and shown on the abstract page https://arxiv.org/abs/2609.04057v1. Full licence notice: \texttt{licenses/arxiv-2609.04057-CC-BY-4.0.txt}.\par\smallskip

\noindent\textit{Editorial scope.} Counted unit: the article's thm thm:general (source0.tex lines 561--582). The article's complete original proof of it is delivered with it (source0.tex lines 1403--1564, one \texttt{proof} environment of 162 lines, which the article places away from the statement). No mathematics is written, repaired or re-derived: the statement and the proof are the article's own lines, in the article's own notation, and no retained line is substituted or re-flowed.  No further source-local context is needed: every cross-reference of every retained line resolves inside the two delivered spans.  Nothing was omitted from any retained span, so the package carries no omission record.  No retained line contains a citation, so the wrapper carries no bibliography and no bibliography entry.  No retained line contains a figure environment, an image-inclusion command or drawing code, so no image is delivered and the package declares no asset dependency.  Editorial formatting only, in addition: an AMS \texttt{amsart} preamble that restates the article's own declarations and macros used by the retained lines, namely the environments \texttt{thm} on separate counters, together with the standard packages those lines need.  The retained environments print in this document as Theorem 1 in order of appearance, whereas the article prints the same items with the numbering of its own full document.  The complete native source of the article as distributed in the arXiv e-print for this exact version is bundled unchanged as \texttt{sources/209340/source0.tex}.
	\begin{thm}
		\label{thm:general}
		Let $\left(\gamma(\beta)\right)_{\beta\geq0}$ be a family of probability vectors of type $\mu=(m_1,\ldots,m_k)$ as $\beta\to\infty$, with multiscale profile $(q^1,\ldots,q^k)$. Then the semigroups $\mathcal{S}^{\gamma(\beta)}$
		converge, in the Kuratowski sense, to the set $\mathcal T^{q^1,\ldots,q^k}$ of row-stochastic block-upper-triangular matrices 
		\begin{equation}
			P=
			\begin{pmatrix}
				P^{11} & P^{12} & \cdots & \cdots & P^{1k}\\
				0 & P^{22} & \cdots & & \vdots\\
				\vdots & \vdots & \ddots & & \vdots\\
				\vdots & & & P^{k-1,k-1} & P^{k-1,k}\\
				0 & \cdots & \cdots & 0 & P^{kk}
			\end{pmatrix},
		\end{equation}
		where each block $P^{ab}$ has size $m_a\times m_b$, and the diagonal blocks $P^{aa}$ are $q^a$-subpreserving:
		\begin{equation}
			q^aP^{aa}\leq q^a
		\end{equation}
		for all $a=1,\ldots,k$. (In particular, the last block $P^{kk}$ is $q^k$-preserving.)
		
		As a consequence,  for $n\geq3$, the set-valued map $q\mapsto \mathcal{S}^q$ is not Painlev\'e-Kuratowski continuous at the boundary points of $\Delta_{n-1}$. 
	\end{thm}

			\begin{proof}[Proof of Theorem~\ref{thm:general}]
				Since $
				\operatorname*{Li}\mathcal{S}^{\gamma(\beta)}
				\subseteq
				\operatorname*{Ls}\mathcal{S}^{\gamma(\beta)}
				$,
				to show that
				$$
				\operatorname*{Lim}\mathcal{S}^{\gamma(\beta)}
				=
				\mathcal T^{q^1,\ldots,q^k}
				$$
				it suffices to prove the two inclusions
				\begin{equation}
					\label{eq:2_inclusions}
					\operatorname*{Ls}\mathcal{S}^{\gamma(\beta)}
					\subseteq
					\mathcal T^{q^1,\ldots,q^k},
					\qquad
					\mathcal T^{q^1,\ldots,q^k}
					\subseteq
					\operatorname*{Li}\mathcal{S}^{\gamma(\beta)}.
				\end{equation}
				We begin with the upper limit. Suppose that
				$$
				P\in\operatorname*{Ls}\mathcal{S}^{\gamma(\beta)}.
				$$
				Then, there exists a sequence $\beta_\ell\to\infty$ and matrices
				$P^{\beta_\ell}\in S^{\gamma(\beta_\ell)}$
				such that $P^{\beta_\ell}\to P$.
				Since
				$$
				\gamma(\beta_\ell)P^{\beta_\ell}
				=
				\gamma(\beta_\ell),
				$$
				for every block index $b=1,\ldots,k$,
				\begin{equation}
					\label{eq:block-balance}
					\sum_{a=1}^k
					\gamma^{a}(\beta_\ell)
					\bigl(P^{\beta_\ell}\bigr)^{ab}
					=
					\gamma^{b}(\beta_\ell),
				\end{equation}
				where $\gamma^{a}(\beta_\ell)$ and 
				$\bigl(P^{\beta_\ell}\bigr)^{ab}$ are conformable for row-by-column multiplication.
				We first show that all strictly lower-triangular blocks of $P^{\beta_{\ell}}$ vanish in the limit.
				
				Fix indices
				$$
				1\le b<a\le k.
				$$
				Since all matrix entries are nonnegative, equation~\eqref{eq:block-balance} implies
				$$
				\gamma^{a}(\beta_\ell)
				\bigl(P^{\beta_\ell}\bigr)^{ab}
				\le
				\gamma^{b}(\beta_\ell),
				$$
				componentwise. Recall that for $b<a$, $$\|\gamma^b(\beta_{\ell})\|_1\ll \|\gamma^a(\beta_{\ell})\|_1,$$ as $\ell\to\infty$. Therefore,
				$$
				\bigl(P^{\beta_\ell}\bigr)^{ab}\to0,
				$$
				and passing to the limit gives
				$$
				P^{ab}=0,
				\quad \text{for $b<a$}.
				$$
				Thus $P$ is block upper triangular.
				
				We now consider the diagonal blocks. Again from~\eqref{eq:block-balance},
				taking $b=a$, we obtain
				$$
				\gamma^{a}(\beta_\ell)
				\bigl(P^{\beta_\ell}\bigr)^{aa}
				\le
				\gamma^{a}(\beta_\ell).
				$$
				Therefore every diagonal block $\bigl(P^{\beta_\ell}\bigr)^{aa}$
				is $\gamma^a(\beta_{\ell})$-subpreserving. Passing to the limit shows that each diagonal block $P^{aa}$ is $q^a$-subpreserving. We conclude that $
				P\in\mathcal T^{q^1,\ldots,q^k}$.
				This proves
				$$
				\operatorname*{Ls}\mathcal{S}^{\gamma(\beta)}
				\subseteq
				\mathcal T^{q^1,\ldots,q^k}.
				$$
				
				We now prove the other inclusion.
				Let $P\in\mathcal T^{q^1,\ldots,q^k}$. 
				We construct matrices
				$P^{\beta}\in \mathcal{S}^{\gamma(\beta)}$
				such that
				$P^{\beta}\to P$.
				The idea is the following: we keep all rows of $P$ except the last one, and modify the final row so as to enforce exact Gibbs invariance.
				Let
				$$
				p=(P_{n1},\ldots,P_{nn})
				$$
				be the last row of $P$, and let
				$$
				p^{(\beta)}=(p^{(\beta)}_1,\ldots,p^{(\beta)}_n)
				$$
				be defined by
				\begin{equation}
					\label{eq:pbeta_def}
					p^{(\beta)}=\frac{1}{\gamma_n(\beta)}\left[\gamma(\beta)-\gamma(\beta)P\right]+p.
				\end{equation}
				In coordinates,
				\begin{equation}
					\label{eq:pbeta}
					p^{(\beta)}_j
					=
					\frac{
						\gamma_j(\beta)
						-
						\sum_{i=1}^{n-1}
						\gamma_i(\beta)P_{ij}
					}
					{\gamma_n(\beta)},
					\qquad
					j=1,\ldots,n.
				\end{equation}
				Define $P^{\beta}$ by replacing the last row $p$ of $P$ with $p^{(\beta)}$:
				\begin{equation}
					P^{\beta}_{ij}
					=
					\begin{cases}
						P_{ij}, & i<n,\\
						p^{(\beta)}_j, & i=n.
					\end{cases}
				\end{equation}
				By construction, $P^{\beta}$ satisfies
				$$
				\gamma(\beta)P^{\beta}
				=
				\gamma(\beta),
				$$
				for all $\beta\geq0$. 
				We claim that $P^{\beta}$ has row sums equal to $1$, and nonnegative entries for all sufficiently large $\beta$. The claim is true for the first $n-1$ rows of $P^{\beta}$. Since $P$ is row-stochastic, from~\eqref{eq:pbeta_def} we have that the last row $p^{(\beta)}$ of $P^{\beta}$ sums to $1$. 
				Positivity for sufficiently large $\beta$ follows from the $q^a$-subpreserving property of the diagonal blocks $P^{aa}$.
				Hence, we exhibited a family $P^{(\beta)}$ of  $\gamma(\beta)$-preserving row-stochastic matrices such that  $P^{(\beta)}\to P$.
				Thus
				$$
				P\in
				\operatorname*{Li}\mathcal{S}^{\gamma(\beta)}.
				$$
				Since $P\in\mathcal T^{q^1,\ldots,q^k}$ was arbitrary,
				$$
				\mathcal T^{q^1,\ldots,q^k}
				\subseteq
				\operatorname*{Li}\mathcal{S}^{\gamma(\beta)}. 
				$$
				Combining the two inclusions~\eqref{eq:2_inclusions}
				we conclude that
				\begin{equation}
					\operatorname{Lim}S^{\gamma(\beta)}= \mathcal T^{q^1,\ldots,q^k}.
				\end{equation}
				
				
			\end{proof}

\end{document}
